\section{Fourier transform}

\subsection{From a series to an integral}

A function on the whole line need not repeat. One can still compare it with the complex Fourier series \eqref{eq:ch05-complex-series} by pretending that it is periodic with a very large period $2\ell$ and then letting $\ell\to\infty$. The discrete frequencies
\begin{equation}
\omega_n=\frac{n\pi}{\ell},
\qquad
\Delta\omega=\omega_n-\omega_{n-1}=\frac{\pi}{\ell},
\end{equation}
become denser as the period grows. The complex coefficient on that period is
\begin{equation}
c_n=\frac{1}{2\ell}\int_{-\ell}^{\ell}f(x)e^{-\mathrm{i}\omega_n x}\,\mathrm{d}x,
\end{equation}
and the series is $f(x)=\sum_n c_n e^{\mathrm{i}\omega_n x}$. The spacing and the normalization are tied together by
\begin{equation}
\frac{1}{2\ell}=\frac{\Delta\omega}{2\pi}.
\end{equation}
Thus
\begin{equation}
f(x)=\sum_{n=-\infty}^{\infty}
\left(\frac{1}{2\pi}\int_{-\ell}^{\ell}f(s)e^{-\mathrm{i}\omega_n s}\,\mathrm{d}s\right)
e^{\mathrm{i}\omega_n x}\,\Delta\omega.
\end{equation}
As $\ell\to\infty$ the sum becomes an integral and the inner integral extends over the line. The resulting asymmetric pair is
\begin{align}
F_{\mathrm{asym}}(\omega)
&=\int_{-\infty}^{\infty}f(x)e^{-\mathrm{i}\omega x}\,\mathrm{d}x,\\
f(x)
&=\frac{1}{2\pi}\int_{-\infty}^{\infty}F_{\mathrm{asym}}(\omega)e^{\mathrm{i}\omega x}\,\mathrm{d}\omega.
\end{align}
The factor $2\pi$ sits on the inverse transform because it came from $\Delta\omega/(2\pi)$, not from the forward coefficient. Placing $2\pi$ on the forward transform, or calling the frequency density the inverse transform, mixes the two sides of this limit.

It is convenient to split the constant symmetrically, writing $F(\omega)=F_{\mathrm{asym}}(\omega)/\sqrt{2\pi}$. The \term{Fourier transform} and its inverse used throughout these notes are then
\begin{align}
F(\omega)&=\frac{1}{\sqrt{2\pi}}\int_{-\infty}^{\infty}f(x)e^{-\mathrm{i}\omega x}\,\mathrm{d}x,
\label{eq:ch06-ft}\\
f(x)&=\frac{1}{\sqrt{2\pi}}\int_{-\infty}^{\infty}F(\omega)e^{\mathrm{i}\omega x}\,\mathrm{d}\omega.
\label{eq:ch06-ift}
\end{align}
One abbreviates $F=\mathcal{F}[f]$ and $f=\mathcal{F}^{-1}[F]$. The function $f$ is the original and $F$ its transform. Differentiation will bring down a factor $\mathrm{i}\omega$ with no extra constant, which is why this sign convention, $e^{-\mathrm{i}\omega x}$ in the forward transform, is kept fixed.

The Fourier integral theorem gives a usable sufficient condition. If $f$ satisfies the Dirichlet conditions on every finite interval and is absolutely integrable on the real line, then the inversion integral reproduces $f(x)$ at each point of continuity, and it reproduces the average of the two one-sided limits at a jump. Absolute integrability excludes, for example, a growing exponential $e^{\alpha x}$ with $\alpha>0$. That function is taken up by the Laplace transform later in the chapter.

There is an equivalent real form. Combining positive and negative frequencies gives
\begin{equation}
f(x)=\int_0^{\infty}A(\omega)\cos(\omega x)\,\mathrm{d}\omega
+\int_0^{\infty}B(\omega)\sin(\omega x)\,\mathrm{d}\omega,
\end{equation}
where
\begin{equation}
A(\omega)=\frac{1}{\pi}\int_{-\infty}^{\infty}f(x)\cos(\omega x)\,\mathrm{d}x,
\qquad
B(\omega)=\frac{1}{\pi}\int_{-\infty}^{\infty}f(x)\sin(\omega x)\,\mathrm{d}x.
\end{equation}
Writing $C\cos(\omega x-\varphi)$ with $C=\sqrt{A^2+B^2}$ and $\varphi=\arctan(B/A)$ separates the \term{amplitude spectrum} $C(\omega)$ from the \term{phase spectrum} $\varphi(\omega)$. An even original has $B=0$ and a cosine integral; an odd original has $A=0$ and a sine integral. An even real function has a real even transform, which is often the quickest check that a calculation has not picked up a spurious phase.

\subsection{Elementary transforms}

The three calculations in this subsection are the ones used repeatedly later: a two-sided exponential, a unit impulse, and a Gaussian.

\begin{example}[Two-sided exponential decay]
Let $f(x)=e^{-\alpha|x|}$ with $\alpha>0$. Split the integral at the origin:
\begin{align}
F(\omega)
&=\frac{1}{\sqrt{2\pi}}
\left(
\int_{-\infty}^{0}e^{\alpha x}e^{-\mathrm{i}\omega x}\,\mathrm{d}x
+\int_0^{\infty}e^{-\alpha x}e^{-\mathrm{i}\omega x}\,\mathrm{d}x
\right)
\nonumber\\
&=\frac{1}{\sqrt{2\pi}}
\left(\frac{1}{\alpha+\mathrm{i}\omega}+\frac{1}{\alpha-\mathrm{i}\omega}\right)
=\sqrt{\frac{1}{2\pi}}\,\frac{2\alpha}{\alpha^2+\omega^2}.
\end{align}
The transform is real and even, as it must be for a real even original. Increasing $\alpha$ makes $f$ more concentrated about $x=0$ and makes $F$ decay more slowly. Localization in $x$ and localization in $\omega$ move in opposite directions.
\end{example}

A unit impulse is not a function, but its action under an integral is unambiguous and is enough to compute the transform. The \term{Dirac delta} satisfies
\begin{equation}
\int_{-\infty}^{\infty}\delta(x)\,\mathrm{d}x=1,
\qquad
\int_{-\infty}^{\infty}f(x)\delta(x-a)\,\mathrm{d}x=f(a)
\end{equation}
for continuous $f$. The second identity is the \term{sifting property}. One also has $\delta(-x)=\delta(x)$ and, for $a\neq 0$,
\begin{equation}
\delta(ax)=\frac{\delta(x)}{|a|}.
\end{equation}
A standard smooth approximation is the Lorentzian family $\varepsilon/[\pi(\varepsilon^2+x^2)]$, which has unit area for every $\varepsilon>0$ and concentrates at the origin as $\varepsilon\to 0$. The next section develops these facts, together with derivatives and composition; the transform itself is only the sifting property.

\begin{example}[Transform of a delta function]
\begin{equation}
\mathcal{F}[\delta(x)]
=\frac{1}{\sqrt{2\pi}}\int_{-\infty}^{\infty}\delta(x)e^{-\mathrm{i}\omega x}\,\mathrm{d}x
=\frac{1}{\sqrt{2\pi}}.
\end{equation}
An impulse localized at a single point has a perfectly flat spectrum. Shifting the impulse multiplies the transform by a phase, $\mathcal{F}[\delta(x-x_0)]=e^{-\mathrm{i}\omega x_0}/\sqrt{2\pi}$, which is the delay rule of the next subsection.
\end{example}

\begin{example}[A Gaussian transforms to a Gaussian]
Let $f(x)=e^{-\alpha x^2}$ with $\alpha>0$. Complete the square in the exponent of \eqref{eq:ch06-ft}:
\begin{equation}
-\alpha x^2-\mathrm{i}\omega x
=-\alpha\left(x+\frac{\mathrm{i}\omega}{2\alpha}\right)^2-\frac{\omega^2}{4\alpha}.
\end{equation}
The contour shift $t=x+\mathrm{i}\omega/(2\alpha)$ does not change the integral of the Gaussian, whose integral on the real line is $\sqrt{\pi/\alpha}$. Therefore
\begin{equation}
F(\omega)
=\frac{1}{\sqrt{2\pi}}e^{-\omega^2/(4\alpha)}\sqrt{\frac{\pi}{\alpha}}
=\frac{1}{\sqrt{2\alpha}}\,e^{-\omega^2/(4\alpha)}.
\label{eq:ch06-gaussian}
\end{equation}
The transform is again a Gaussian. The product of the decay constants in the two exponents is fixed: if the original narrows, the transform widens, and conversely. This is the same tradeoff seen for the exponential, now with an explicit width $\alpha^{-1/2}$ in $x$ and width $\alpha^{1/2}$ in $\omega$.
\end{example}

\begin{example}[A rectangular pulse]
Let $f(x)=1$ for $|x|\leq a$ and $f(x)=0$ otherwise, with $a>0$. Then
\begin{equation}
F(\omega)
=\frac{1}{\sqrt{2\pi}}\int_{-a}^{a}e^{-\mathrm{i}\omega x}\,\mathrm{d}x
=\sqrt{\frac{2}{\pi}}\,\frac{\sin(a\omega)}{\omega}
=\sqrt{\frac{2}{\pi}}\,a\,\frac{\sin(a\omega)}{a\omega}.
\end{equation}
At $\omega=0$ the quotient tends to $\sqrt{2/\pi}\,a$, which is also $F(0)=(1/\sqrt{2\pi})\int f$. The transform oscillates and decays as $1/|\omega|$. A wider pulse has a narrower central lobe. This one-dimensional pattern is the model for the rectangular aperture computed below.
\end{example}

\subsection{Basic properties}

Assume $\mathcal{F}[f]=F$ and that $f$ and its derivatives decay fast enough for the boundary terms written below to vanish.

Differentiation under the inversion integral \eqref{eq:ch06-ift}, or an integration by parts in \eqref{eq:ch06-ft}, gives the derivative theorem
\begin{equation}
\mathcal{F}[f']=\mathrm{i}\omega F(\omega).
\label{eq:ch06-derivative}
\end{equation}
The boundary term $\bigl[f(x)e^{-\mathrm{i}\omega x}\bigr]_{-\infty}^{\infty}$ is zero by the decay assumption, and the remaining integral is $\mathrm{i}\omega$ times $F(\omega)$. Each extra derivative multiplies by another factor $\mathrm{i}\omega$. This is the reason a linear differential equation with constant coefficients becomes an algebraic equation for $F$.

If $F(0)=0$ and the running integral $g(x)=\int_{-\infty}^{x}f(s)\,\mathrm{d}s$ tends to zero at infinity, the derivative theorem applied to $g$ yields
\begin{equation}
\mathcal{F}[g]=\frac{F(\omega)}{\mathrm{i}\omega}.
\end{equation}
The hypothesis $F(0)=0$ says that $f$ has integral zero, so $g(\infty)=0$. Without it, the running integral tends to a nonzero constant and is not a classical Fourier-transformable function.

The remaining rules are changes of the variable or multiplications by exponentials. For $a\neq 0$,
\begin{equation}
\mathcal{F}[f(ax)]=\frac{1}{|a|}F\Bigl(\frac{\omega}{a}\Bigr).
\label{eq:ch06-similarity}
\end{equation}
The absolute value is the Jacobian. A negative $a$ reverses the axis and would be missed by writing $1/a$. For a shift $x_0$,
\begin{equation}
\mathcal{F}[f(x-x_0)]=e^{-\mathrm{i}\omega x_0}F(\omega).
\label{eq:ch06-delay}
\end{equation}
The proof is the substitution $u=x-x_0$ in \eqref{eq:ch06-ft}: the Jacobian is $1$, and $e^{-\mathrm{i}\omega x}=e^{-\mathrm{i}\omega(u+x_0)}$ factors out of the integral. Modulation is the companion statement, obtained by writing the extra exponential into the kernel,
\begin{equation}
\mathcal{F}[e^{\mathrm{i}\omega_0 x}f(x)]=F(\omega-\omega_0).
\end{equation}
A shift in time is a phase in frequency; a modulation in time is a shift in frequency.

\begin{example}[The delay rule applied to a pulse]
The rectangular pulse of width $2a$ centered at $x_0$ is the centered pulse shifted by $x_0$. Its transform is therefore
\begin{equation}
e^{-\mathrm{i}\omega x_0}\sqrt{\frac{2}{\pi}}\,\frac{\sin(a\omega)}{\omega}.
\end{equation}
The amplitude spectrum is unchanged. Only the phase $-\omega x_0$ records where the pulse sits. This is why a pure delay is invisible to an intensity measurement $|F(\omega)|^2$.
\end{example}

The \term{convolution} used here is
\begin{equation}
(f_1*f_2)(x)=\int_{-\infty}^{\infty}f_1(\alpha)f_2(x-\alpha)\,\mathrm{d}\alpha.
\label{eq:ch06-convolution}
\end{equation}
The arguments are $x-\alpha$ and $\alpha$, not $\alpha-x$. The reflected convention $\int f_1(\alpha)f_2(\alpha-x)\,\mathrm{d}\alpha$ convolves $f_2$ with a reflected copy of $f_1$ and does not obey the product rule below. Substitute the inversion formula for $f_2$ and change the order of integration:
\begin{align}
(f_1*f_2)(x)
&=\int f_1(\alpha)
\left(\frac{1}{\sqrt{2\pi}}\int F_2(\omega)e^{\mathrm{i}\omega(x-\alpha)}\,\mathrm{d}\omega\right)\mathrm{d}\alpha
\nonumber\\
&=\frac{1}{\sqrt{2\pi}}\int F_2(\omega)e^{\mathrm{i}\omega x}
\left(\int f_1(\alpha)e^{-\mathrm{i}\omega\alpha}\,\mathrm{d}\alpha\right)\mathrm{d}\omega
\nonumber\\
&=\frac{1}{\sqrt{2\pi}}\int\bigl(\sqrt{2\pi}\,F_1(\omega)F_2(\omega)\bigr)e^{\mathrm{i}\omega x}\,\mathrm{d}\omega.
\end{align}
The inner integral is $\sqrt{2\pi}\,F_1(\omega)$. Comparing with \eqref{eq:ch06-ift} gives the convolution theorem for this normalization,
\begin{equation}
\mathcal{F}[f_1*f_2]=\sqrt{2\pi}\,F_1(\omega)F_2(\omega).
\label{eq:ch06-conv-thm}
\end{equation}
The factor $\sqrt{2\pi}$ belongs to the symmetric convention. It is absent from the asymmetric convention in which $2\pi$ sits entirely on one side, and it is absent from the Laplace convolution on a half-line. The next section returns to convolution in order to treat Gaussians, step functions, and products of delta functions; the identity \eqref{eq:ch06-conv-thm} is the form required there.

\subsection{Higher-dimensional transforms}

On $\mathbb{R}^d$ the same symmetric splitting of $(2\pi)^d$ puts the factor $(2\pi)^{-d/2}$ on each side:
\begin{align}
F(\mathbf{k})
&=\frac{1}{(2\pi)^{d/2}}
\int_{\mathbb{R}^d}f(\mathbf{r})e^{-\mathrm{i}\mathbf{k}\cdot\mathbf{r}}\,\mathrm{d}^d\mathbf{r},
\label{eq:ch06-ft-d}\\
f(\mathbf{r})
&=\frac{1}{(2\pi)^{d/2}}
\int_{\mathbb{R}^d}F(\mathbf{k})e^{\mathrm{i}\mathbf{k}\cdot\mathbf{r}}\,\mathrm{d}^d\mathbf{k}.
\end{align}
An equivalent convention places $(2\pi)^{-d}$ on one side only. The two conventions must not be mixed inside a single calculation. In two dimensions the symmetric factor is $1/(2\pi)$ on each side. In three dimensions it is $(2\pi)^{-3/2}$ on each side.

\begin{example}[Fraunhofer pattern of a rectangular aperture]
Let $f(x,y)=A$ on the rectangle $|x|\leq X$, $|y|\leq Y$, and let $f=0$ outside. This is a standard model of a uniformly illuminated rectangular opening. With the two-dimensional symmetric transform,
\begin{align}
F(k_1,k_2)
&=\frac{A}{2\pi}
\int_{-X}^{X}e^{-\mathrm{i}k_1 x}\,\mathrm{d}x
\int_{-Y}^{Y}e^{-\mathrm{i}k_2 y}\,\mathrm{d}y
\nonumber\\
&=\frac{A}{2\pi}
\cdot\frac{2\sin(k_1 X)}{k_1}
\cdot\frac{2\sin(k_2 Y)}{k_2}
\nonumber\\
&=\frac{2AXY}{\pi}
\,\frac{\sin(k_1 X)}{k_1 X}
\,\frac{\sin(k_2 Y)}{k_2 Y}.
\end{align}
Each factor $\sin(u)/u$ is the one-dimensional pulse transform. The pattern is bright on axis, where both quotients equal $1$ and $F(0,0)=2AXY/\pi$, and it has nodal lines $k_1 X=n\pi$ and $k_2 Y=m\pi$ for nonzero integers $n,m$. Using $(2\pi)^{-2}$ in front of this integral would disagree with \eqref{eq:ch06-ft-d} and would not reproduce the algebra above. The asymmetric two-dimensional convention, with $1/(2\pi)$ on the forward transform and no constant on the inverse, is a different pair and is not the one used here.
\end{example}
